\subsubsection{Data Processing}
\begin{enumerate}
    \item \textbf{Normalization of Expression Values}
          
          The subroutine runs a complete normalization pipeline on gene expression data.
          It first applies standard deviation normalization, then quantile normalization,
          followed by replicate averaging, and finally a log2 transformation. The
          result is returned in the \texttt{buf\_log} array, which contains the
          fully normalized expression values. Additionally intermediate results of the
          performed calculations are available in output buffer arrays. Call the
          \texttt{normalization\_pipeline} subroutine with the following arguments:
          
          \textbf{Input arguments:}
          \begin{itemize}
            \item \texttt{n\_genes integer(int32)}: Total number of genes
                  
            \item \texttt{n\_tissues integer(int32)}: Number of axes (tissues /
                  dimensions)
                  
            \item \texttt{input\_matrix real(real64)}: Expression vector matrix with
                  size \texttt{n\_genes $\times$ n\_tissues}
                  
            \item \texttt{max\_stack integer(int32)}: Stack size for quicksort
                  
            \item \texttt{group\_s integer(int32)}: Start column index for each replicate
                  group (\texttt{n\_grps})
                  
            \item \texttt{group\_c integer(int32)}: Number of columns per replicate group
                  (\texttt{n\_grps})
                  
            \item \texttt{n\_grps integer(int32)}: Number of replicate groups
          \end{itemize}
          
          \textbf{Output arguments:}
          \begin{itemize}
            \item \texttt{buf\_stddev real(real64)}: Buffer for std dev normalization
                  (\texttt{n\_genes} * \texttt{n\_tissues})
                  
            \item \texttt{buf\_quant real(real64)}: Buffer for quantile normalization
                  (\texttt{n\_genes} * \texttt{n\_tissues})
                  
            \item \texttt{buf\_avg real(real64)}: Buffer for replicate averaging (\texttt{n\_genes}
                  * \texttt{n\_grps})
                  
            \item \texttt{buf\_log real(real64)}: Buffer for log2(x+1)
                  transformation (\texttt{n\_genes} * \texttt{n\_grps})
                  
            \item \texttt{temp\_col real(real64)}: Temporary column vector for sorting
                  (\texttt{n\_genes})
                  
            \item \texttt{rank\_means real(real64)}: Buffer for rank means (\texttt{n\_genes})
                  
            \item \texttt{perm integer(int32)}: Permutation vector for sorting (\texttt{n\_genes})
                  
            \item \texttt{stack\_left integer(int32)}: Left stack for quicksort (\texttt{max\_stack})
                  
            \item \texttt{stack\_right integer(int32)}: Right stack for quicksort (\texttt{max\_stack})
                  
            \item \texttt{ierr integer(int32)}: Error code (for error details please
                  check \hyperref[sec:error-handling]{error handling section} )
          \end{itemize}
          
          \begin{lstlisting}[language=Fortran, caption={Normalization Pipeline}]
call normalization_pipeline(n_genes, n_tissues, input_matrix, &
                            buf_stddev, buf_quant, buf_avg, buf_log, &
                            temp_col, rank_means, perm, stack_left, &
                            stack_right, max_stack, group_s, &
                            group_c, n_grps, ierr)
          \end{lstlisting}
          
          \textbf{Alternatively you can calculate each step of the normalization
          pipeline manually by calling the following subroutines in sequence:}
          \begin{enumerate}
            \item \textbf{Normalize by Standard Deviation}
                  
                  The subroutine normalizes each gene's expression vector using \texttt{sqrt(mean(x\textasciicircum
                    2))} across tissues (not classical standard deviation). Call the \texttt{normalize\_by\_std\_dev}
                  subroutine with the following arguments:
                  
                  \textbf{Input arguments:}
                  \begin{itemize}
                    \item \texttt{n\_genes integer(int32)}: Number of genes (rows)
                          
                    \item \texttt{n\_tissues integer(int32)}: Number of tissues (columns)
                          
                    \item \texttt{input\_matrix real(real64)}: Flattened input matrix of
                          gene expression values (column-major)
                  \end{itemize}
                  \textbf{Output arguments:}
                  \begin{itemize}
                    \item \texttt{output\_matrix real(real64)}: Output normalized matrix
                          (same shape as input)
                          
                    \item \texttt{ierr integer(int32)}: Error code (for error details
                          please check \hyperref[sec:error-handling]{error handling section}
                          )
                  \end{itemize}
                  \begin{lstlisting}[language=Fortran, caption={Standard Deviation Normalization}]
call normalize_by_std_dev(n_genes, n_tissues, input_matrix, output_matrix, ierr)
                  \end{lstlisting}
                  
            \item \textbf{Quantile Normalization}
                  
                  The subroutine performs quantile normalization of a gene expression matrix
                  (F42-compliant) and computes average expression per rank across tissues.
                  Call the \texttt{quantile\_normalization} subroutine with the following
                  arguments:
                  
                  \textbf{Input arguments:}
                  \begin{itemize}
                    \item \texttt{n\_genes integer(int32)}: Number of genes (rows)
                          
                    \item \texttt{n\_tissues integer(int32)}: Number of tissues (columns)
                          
                    \item \texttt{input\_matrix real(real64)}: Flattened input matrix (column-major)
                          
                    \item \texttt{max\_stack integer(int32)}: Stack size passed from R
                  \end{itemize}
                  \textbf{Output arguments:}
                  \begin{itemize}
                    \item \texttt{output\_matrix real(real64)}: Output normalized matrix
                          (same shape as input)
                          
                    \item \texttt{temp\_col real(real64)}: Temporary vector for column sorting
                          (size \texttt{n\_genes})
                          
                    \item \texttt{rank\_means real(real64)}: Preallocated vector to store
                          rank means (size \texttt{n\_genes})
                          
                    \item \texttt{perm integer(int32)}: Permutation vector (size \texttt{n\_genes})
                          
                    \item \texttt{stack\_left integer(int32)}: Manual quicksort stack ($\geq
                          log2(\texttt{n\_genes}) + 10$)
                          
                    \item \texttt{stack\_right integer(int32)}: Manual quicksort stack (same
                          size as stack\_left)
                          
                    \item \texttt{ierr integer(int32)}: Error code (for error details
                          please check \hyperref[sec:error-handling]{error handling section}
                          )
                  \end{itemize}
                  \begin{lstlisting}[language=Fortran, caption={Quantile Normalization}]
call quantile_normalization(n_genes, n_tissues, input_matrix, output_matrix, temp_col, rank_means, perm, stack_left, stack_right, max_stack, ierr)
                  \end{lstlisting}
                  
            \item \textbf{Calculate Tissue Averages}
                  
                  The subroutine calculates tissue averages by averaging replicates within
                  each group. For each group of tissue replicates, it computes the average
                  expression per gene. The input matrix is column-major, flattened as a 1D
                  array. Call the \texttt{calc\_tiss\_avg} subroutine with the following
                  arguments:
                  
                  \textbf{Input arguments:}
                  \begin{itemize}
                    \item \texttt{n\_gene integer(int32)}: Number of genes (rows)
                          
                    \item \texttt{n\_grps integer(int32)}: Number of tissue groups
                          
                    \item \texttt{group\_s integer(int32)}: Start column index for each group
                          (length: \texttt{n\_grps})
                          
                    \item \texttt{group\_c integer(int32)}: Number of columns per group (length:
                          \texttt{n\_grps})
                          
                    \item \texttt{input\_matrix real(real64)}: Flattened input matrix (length:
                          \texttt{n\_gene} * sum(\texttt{group\_c}))
                  \end{itemize}
                  \textbf{Output arguments:}
                  \begin{itemize}
                    \item \texttt{output\_matrix real(real64)}: Flattened output matrix (length:
                          \texttt{n\_gene} * \texttt{n\_grps})
                          
                    \item \texttt{ierr integer(int32)}: Error code (for error details
                          please check \hyperref[sec:error-handling]{error handling section}
                          )
                  \end{itemize}
                  \begin{lstlisting}[language=Fortran, caption={Tissue Averages Calculation}]
call calc_tiss_avg(n_gene, n_grps, group_s, group_c, input_matrix, &
                    output_matrix, ierr)
                  \end{lstlisting}
                  
            \item \textbf{Log2 Transformation}
                  
                  The subroutine applies \texttt{log2(x + 1)} transformation to each
                  element of the input matrix. This is computed via \texttt{log(x + 1) /
                    log(2)} for compatibility with WebAssembly (WASM). Call the \texttt{log2\_transformation}
                  subroutine with the following arguments:
                  
                  \textbf{Input arguments:}
                  \begin{itemize}
                    \item \texttt{n\_genes integer(int32)}: Number of genes (rows)
                          
                    \item \texttt{n\_tissues integer(int32)}: Number of tissues (columns)
                          
                    \item \texttt{input\_matrix real(real64)}: Flattened input matrix (size:
                          \texttt{n\_genes} * \texttt{n\_tissues})
                  \end{itemize}
                  \textbf{Output arguments:}
                  \begin{itemize}
                    \item \texttt{output\_matrix real(real64)}: Output matrix (same size
                          as input)
                          
                    \item \texttt{ierr integer(int32)}: Error code (for error details
                          please check \hyperref[sec:error-handling]{error handling section}
                          )
                  \end{itemize}
                  \begin{lstlisting}[language=Fortran, caption={Log2 Transformation}]
call log2_transformation(n_genes, n_tissues, input_matrix, &
                        output_matrix, ierr)
                  \end{lstlisting}
          \end{enumerate}
          
    \item \textbf{Calculate Fold Change}
          
          If fold change calculation is needed, this subroutine can be used after
          normalization. The subroutine calculates \texttt{log2} fold changes
          between condition and control columns. For each control-condition pair,
          this subroutine computes the \texttt{log2} fold change by subtracting the expression
          value in the control column from the corresponding value in the condition column,
          for all genes. Call the \texttt{calc\_fchange} subroutine with the
          following arguments:
          
          \textbf{Input arguments:}
          \begin{itemize}
            \item \texttt{n\_genes integer(int32)}: Total number of genes
                  
            \item \texttt{n\_cols integer(int32)}: Number of columns in the input matrix
                  
            \item \texttt{n\_pairs integer(int32)}: Number of control-condition pairs
                  
            \item \texttt{control\_cols integer(int32)}: Control column indices (length
                  \texttt{n\_pairs})
                  
            \item \texttt{cond\_cols integer(int32)}: Condition column indices (length
                  \texttt{n\_pairs})
                  
            \item \texttt{i\_matrix real(real64)}: Input matrix, flattened (length:
                  \texttt{n\_genes} $\times$ \texttt{n\_cols})
          \end{itemize}
          
          \textbf{Output arguments:}
          \begin{itemize}
            \item \texttt{o\_matrix real(real64)}: Output matrix for fold changes (length:
                  \texttt{n\_genes} $\times$ \texttt{n\_pairs})
                  
            \item \texttt{ierr integer(int32)}: Error code (for error details please
                  check \hyperref[sec:error-handling]{error handling section} )
          \end{itemize}
          
          \begin{lstlisting}[language=Fortran, caption={Fold Change Calculation}]
call calc_fchange(n_genes, n_cols, n_pairs, control_cols, cond_cols, &
                    i_matrix, o_matrix, ierr)
          \end{lstlisting}
\end{enumerate}